
relation est équivalente à la relation :
@@M0@@
Or on a :
@@M1@@
D'autre part :
@@M2@@
Comme @@M3@@, on peut appliquer la formule (V (5.4.2)) :
@@M4@@
En posant : @@M5@@, on est donc ramené à montrer la relation :
@@M6@@
c'est-à-dire la relation analogue pour l'immersion @@M7@@, et l'élément @@M8@@. Or par hypothèse @@M9@@ est multiple de @@M10@@; on peut trouver @@M11@@ tel que @@M12@@. Alors @@M13@@ avec @@M14@@ (2.7) ; comme @@M15@@, on a finalement à vérifier :
@@M16@@
La suite exacte
@@M17@@
donne la relation :
@@M18@@
La relation à vérifier s'écrit donc finalement, en appliquant la formule de
projection :
